\documentclass[11pt]{article}

\usepackage[T1]{fontenc}
\usepackage[margin=1in]{geometry}
\usepackage{enumitem}
\usepackage{hyperref}
\usepackage{listings}

\lstset{
  basicstyle=\ttfamily\small,
  breaklines=true,
  columns=fullflexible
}

\title{HW3 Taint Analysis Report}
\author{Jason Chia-Sheng Lin\\Student ID: 513559004}
\date{May 2026}

\begin{document}
\maketitle

\section*{Run Summary}

The completed analyzer emulates the provided \texttt{vuln} binary with Triton,
marks the first 16 bytes of \texttt{user\_input} as tainted, propagates taint
through the implemented \texttt{strncpy} hook, and checks the first 16 bytes of
\texttt{output\_buf}. The observed result in \texttt{output.txt} is:

\begin{lstlisting}
Result: 16 / 16 bytes tainted at sink
\end{lstlisting}

This means every tracked byte that started from the attacker-controlled source
reached the sink.

\section*{Question 1: Taint Propagation Path}

The program's data path has five stages:

\begin{lstlisting}
user_input -> process_data() -> processed -> strncpy -> output_buf
\end{lstlisting}

The array \texttt{user\_input} is the taint source, so the analyzer marks its
first 16 bytes as attacker-controlled before emulation starts. During
\texttt{process\_data()}, each byte of \texttt{processed} is computed from the
corresponding byte of \texttt{user\_input}:

\begin{lstlisting}
processed[i] = user_input[i] ^ 0x42;
\end{lstlisting}

The XOR operation changes the concrete byte value, but the result still depends
directly on attacker-controlled input. XOR with a constant is a deterministic
transformation, not a validation or sanitization step. If an attacker controls
\texttt{user\_input[i]}, then the attacker also influences
\texttt{processed[i]}, even though the byte value is transformed. For that
reason, taint should propagate from \texttt{user\_input[i]} to
\texttt{processed[i]}.

After \texttt{process\_data()}, \texttt{copy\_to\_output()} calls
\texttt{strncpy(output\_buf, processed, 16)}. This copies the transformed bytes
from \texttt{processed} into \texttt{output\_buf}. Because those source bytes are
tainted, the corresponding sink bytes in \texttt{output\_buf} should also be
tainted. The final output confirms that all 16 tracked sink bytes are tainted.

\section*{Question 2: Why the \texttt{strncpy} Hook Propagates Taint Manually}

Triton automatically propagates taint through instructions that it emulates.
However, this assignment does not load or emulate the real libc implementation
of \texttt{strncpy}. Instead, the script redirects that library call to a Python
hook. The hook changes Triton's concrete memory state directly, outside the
normal emulated instruction stream.

Because of that boundary, the hook must copy both the concrete byte value and
the taint state for each byte. In my implementation, each \texttt{src + i} byte
is read, written to \texttt{dest + i}, and then the destination byte is tainted
or untainted to match the source byte.

If the hook copied only concrete bytes and did not copy taint, the analyzer
would produce a false negative. The program would still move attacker-controlled
data into \texttt{output\_buf}, but Triton would not know that the Python hook
preserved the data dependency. The sink could incorrectly appear clean.

\section*{Question 3: Real-World Vulnerability Scenario}

A common real-world scenario is SQL injection in a web application. The taint
source could be an HTTP request parameter such as \texttt{name} or
\texttt{user\_id}. The application might URL-decode the parameter, trim
whitespace, concatenate it into a SQL string, and then pass the final string to
a database execution function.

In this case, the taint sink is the database query execution call, such as
\texttt{execute(query)}. The expected propagation path is:

\begin{lstlisting}
HTTP request parameter -> parsing/decoding -> string concatenation -> SQL query -> execute()
\end{lstlisting}

Dynamic taint analysis can help by showing whether attacker-controlled bytes
from the request reach the query execution sink without being safely handled. If
tainted bytes are present in SQL syntax positions, such as operators, quotes, or
conditions, the result is a strong signal that the code needs parameterized
queries or stricter input handling. Taint analysis does not by itself prove
every exploit, but it points reviewers to dangerous source-to-sink flows.

\section*{Question 4: Limitations and Complementary Techniques}

First, dynamic taint analysis only observes paths that actually execute. If a
test input does not reach a vulnerable branch, the analyzer may miss a real
problem. Fuzzing, unit tests with security assertions, and symbolic execution
can improve path coverage by generating more inputs and exercising more control
flow.

Second, runtime overhead and environment complexity can be high. Instruction
emulation, taint metadata, and library hooks can make analysis slower than
normal execution, and missing hooks can create false negatives or inaccurate
results. Static analysis, code review, and focused integration tests can help
cover library behavior and larger application structure without emulating every
instruction.

Third, implicit flows are difficult. Data may influence control decisions rather
than being copied directly, such as a tainted condition choosing which clean
constant to write. Basic byte-level taint tracking can miss or over-approximate
these cases. Complementary manual review and policy-aware static analysis are
useful for deciding whether such control dependencies matter for the security
property being checked.

\end{document}
